\subsection{Propagation across the components of the rank-four space}
\label{subsec:rh-four-component-propagation}

We now prove the component-support step for the original fitted deep
tower $T$. Let $F=T.F$, $d=[F:\mathbb Q]$, $r=\#T.\iota$, and
$d_v=d/r$ (the common local degree at the ordinary places).
Keep the original complete coefficient DVR $\mathcal O$ and
residue field $k$, the original frame $\Phi=T.\mathrm{fr}(\varnothing)$,
and $X$ with its coefficient ring $B=X.B$ and entry ring $A=X.A$.
Let $J$ be the sum ideal of Theorem~\ref{thm:rh-four-image-ideal}.
We assume precisely the original \texttt{Fr.Bounds} with constant
$C$ and the threshold $d\geq30(C+1)^2$. In particular $r\leq C$.
These are hypotheses on the original tower, not bounds obtained
by changing the deformation space.

\begin{lemma}[The actual completed chart]
\label{lem:rh-four-chart-deficit}
At any original residue evaluation $x:A\to k$ satisfying
\texttt{IsRes}, write $R=\widehat{A_{\ker x}}$.
There is a presentation
\[
 R=\mathcal O[[u_1,\ldots,u_a]]/(f_1,\ldots,f_b),
 \qquad a-b\geq 2+2d-2r.
\]
The map $A\to R$ is flat, and $R$ represents the original lifts
of the residue pair at $x$.
\end{lemma}
\begin{proof}
By the original space contract, $A$ is a finite type algebra over
the complete Noetherian local ring $B$, and its representation
and flag entries generate it over $B$. Two maps agreeing on
these entries agree on $B$ by \texttt{genB} (the characteristic
polynomials and the labeled characters are determined by the
entries). Hence the original lifting property gives uniqueness
as well as existence. A map inducing $x$ on the residue field
extends uniquely to the displayed local completion, for any
complete local target with that residue field. Localization and
completion are flat. This identifies the chart with the actual
marked deformation functor, including the local flags; it is
not a presentation of the semisimple character space.

Here is the relation count. First omit the local \texttt{Mu}
equations. Write $M=\mathfrak{sp}_4(k)$ with the adjoint action
of the residual pair. At an ordinary place let $\mathfrak b_v$
be the Borel Lie algebra defined by its residual flag.
The dimensions of $M$ and $\mathfrak b_v$ are ten and six;
the latter assertion follows from the four positive root
spaces and the two-dimensional diagonal torus. We use the
same relative obstruction argument as for the earlier global
deformation presentation, with the following local complexes.
At an ordinary place the quotient of local cochains is
$C^\bullet(G_v,M/\mathfrak b_v)$. Away from $\ell$ it is the
quotient of $C^\bullet(G_v,M)$ by the inflated cochains of
$G_v/N_v$, where $N_v=I_v\cap\ker \bar r_{\mathrm{diag}}$ is
the original killed subgroup. The action of $N_v$ is trivial.
In the fitted deep tower the original diagonal characters
on this tame inertia quotient have prime-to-$\ell$ order:
the cyclotomic character is unramified there and the finite
characters $\xi,\psi|_{G_F}$ have prime-to-$\ell$ order.
Consequently $G_v/N_v$ is an extension of a procyclic
unramified group by a finite group of order prime to $\ell$.
Its $\ell$-cohomological dimension is at most one and its
Euler characteristic with coefficients in $M$ is zero.
Use the shifted cone of restriction from global cochains
(unramified outside the original finite ramification set,
including $\ell$) to these local quotients.

The residual representation is triangular after the fixed
reordering appropriate to \texttt{Res.shape}. Hence the
composition factors of its matrix adjoint are the quotients
of the four diagonal characters. Since $\ell\ne2$, the
symplectic adjoint is a direct summand of the matrix adjoint,
using the symplectic transpose. The hypothesis \texttt{Ok.g2}
excludes the residual cyclotomic character among all these
quotients (and among the trivial factors). Thus
$H^0(F,M^*(1))=0$. Poitou--Tate duality makes restriction
in degree two onto the full finite local cohomology
surjective. Passing to the displayed quotients preserves
surjectivity, by local cohomological dimension two. The
shifted cone has no degree-three cohomology.

At a real place complex conjugation has multiplier $-1$,
so its two eigenspaces are complementary Lagrangians of
dimension two. Its adjoint fixed space in $M$ has dimension
four, being the endomorphisms of one Lagrangian with the
opposite action on the other. Its negative eigenspace has
dimension six. The global Euler characteristic is therefore
$-6d$. The local quotients away from $\ell$ have Euler
characteristic zero, whereas the ordinary local quotients
have Euler characteristic $-4d_v$. If $h^i$ are the cone
cohomology dimensions, subtraction gives
\[
                       h^1-h^2-h^0=2d.                 \tag{CP1}
\]
All these groups are finite-dimensional by finite-module
Galois cohomology. The same argument works if the tame
local conditions are first imposed on the maximal allowed
ramification quotient.

For clarity, the obstruction and tangent interpretations here
retain the flags. Over a small extension lift the global
matrices into the symplectic similitude group with the fixed
multiplier, possible pointwise by smoothness. Their
multiplication defect is a two-cocycle with coefficients in
$M$. Lift the local flags in the smooth flag chart and
compare the defect with the local Borel, producing its
degree-one primitive in $M/\mathfrak b_v$. At an old tame
place choose a lift on the inertia quotient (its obstruction
vanishes in degree two), and compare with the lifted global
matrices in the quotient of cochains just described.
These form the relative two-cocycle. Changing the lifts
changes it by a relative coboundary. Vanishing permits
correction of the matrices and flags and gives precisely a
lift satisfying all the local conditions. In degree one this
same construction gives the tangent space of pairs, of
dimension $h^1+\dim M-h^0$ before marking, because quotienting
by infinitesimal global conjugation subtracts the dimension
of the coboundaries. The standard small-extension
relation-functional argument bounds the number of minimal
relations by $h^2$: push out the kernel of a minimal
power-series presentation along a linear functional and
observe that vanishing of its obstruction gives a section,
forcing that functional to vanish by minimality. This is
the integral argument of the earlier relative presentation,
not a calculation on its reduced quotient.

The mark is a diagonal matrix with distinct residual
eigenvalues. Conjugation by the eight off-diagonal root
directions maps bijectively to the eight off-diagonal
marking conditions (the factors are differences of distinct
eigenvalues). It therefore both eliminates eight tangent
variables and normalizes any lift across a small extension,
without increasing the obstruction space. The marked
deficit before the \texttt{Mu} equations is at least
$h^1-h^2-h^0+2=2d+2$.

Finally the abelianized local decomposition group, by local
reciprocity, has cyclic $\ell$-primary torsion: it is the
$\ell$-power roots of unity in the local field. The set
\texttt{Mu} is its inverse image (an element whose
$\ell$-power lies in the closed commutator). All diagonal
characters kill that commutator by continuity. The
multiplier is trivial on this torsion by the deep condition
and the absence of $\ell$-torsion in the cyclotomic
principal units. At most two equations, saying that the
first two diagonal characters kill a lift of the torsion
generator, suffice at each place; the other two follow
from the paired product relation. Subtracting at most $2r$
from the deficit proves the lemma.
\end{proof}

\begin{lemma}[Reducible loci in the completed chart]
\label{lem:rh-four-red-dimension}
Suppose in addition that $x$ satisfies \texttt{Good}. In the
chart of the preceding lemma, the reducible locus of the
generic rank-four representation is contained in a union of
nine closed loci, each of dimension at most $d+20C+1$.
\end{lemma}
\begin{proof}
Use zero-based indices and the weights
\[
 (1,1),(0,1),(1,0),(0,0)
\]
on the four coordinates. Write $A_0=(1,0)$, $P=(0,1)$,
$B_0=(-1,1)$, $H=(1,1)$ for the four source root names
$wAl,wPs,wBe,wTh$. The eight roots are these and their
negatives. The marked matrix remains diagonal with distinct
eigenvalues over every prime quotient of the local ring.
An invariant subspace over an algebraic closure of a fraction
field must therefore be spanned by a nonempty proper subset
$I$ of coordinate vectors. Its symplectic orthogonal is the
coordinate subset complementary to $\{3-i:i\in I\}$, and is
invariant too. For a root set $S$ let the corresponding
closed locus impose all entries outside $S$ off the diagonal.
The two invariance conditions give the following list of
possible supports, writing a minus sign for the opposite
root:
\[
\begin{array}{c|l|l}
 & S & I\text{ (up to replacing by its orthogonal)}\\ \hline
1&B_0,-B_0,-A_0,-H,-P&3\\
2&B_0,-B_0,A_0,H,P&0\\
3&H,-H,P,B_0,-A_0&1\\
4&H,-H,-P,-B_0,A_0&2\\
5&A_0,-A_0,-H,-P,-B_0&23\\
6&A_0,-A_0,H,P,B_0&01\\
7&P,-P,B_0,-A_0,-H&13\\
8&P,-P,-B_0,A_0,H&02\\
9&B_0,-B_0,H,-H&03,\ 12
\end{array}                                                     \tag{CP2}
\]
Indeed delete a root as soon as one entry of that weight is
forced zero by either invariant coordinate subspace. This
gives exactly the nine \texttt{redTypes} in the source.

To bound a locus $S$, it suffices to bound the relative
tangent dimension at the residue pair $c$ by $d+20C$;
the coefficient uniformizer adds at most one to the
embedding dimension. If the locus does not contain $c$,
its local quotient is zero. Otherwise write a tangent
as $Z(g),Y_v$ in $M$, respectively the cocycle and the
local flag variation. Its exact conditions are: $Z$
is a locally constant cocycle killing the prescribed
inertia and diagonal at the mark; $Z(g)c.r(g)$ has
support $S$; $Y_v$ belongs to the conjugate strictly
lower Lie algebra in the flag basis; and
\[
 c.n(v)^{-1}\bigl[Z(s)+\operatorname{Ad}(c.r(s))Y_v-Y_v\bigr]c.n(v)
 \quad\text{is upper triangular}\quad(s\in G_v).       \tag{CP3}
\]
We may drop the infinitesimal \texttt{Mu} constraints when
bounding this tangent.

There is a useful elementary degeneration. Conjugate by
$\operatorname{diag}(t^{m_0},\ldots,t^{m_3})$, with
$m_i+m_{3-i}$ constant and every nonzero entry of the
residual matrices and flags of nonnegative degree.
Take their degree-zero parts. Filter the tangent space by
the lowest degree of its entries, including the flag
variations. For an element of lowest degree $j$, the
degree-$j$ part satisfies exactly the tangent equations
at the limiting residual pair. This follows by taking
the lowest homogeneous degree in the cocycle and
(CP3); the equations involving inverses use the same
nonnegative filtration for the inverse matrices, and
the symplectic transpose preserves degree. Locally
constant functions remain locally constant under the
coordinate projections. The successive filtered
quotients inject into the distinct homogeneous subspaces
of the limiting tangent. Thus its dimension bounds
the original tangent dimension, even without choosing
finitely many generators of the Galois group.

In the first two \texttt{Good} alternatives the support
$S$ contains either $-P$ and one of $B_0,-A_0$, or
both $-A_0,-B_0$. For residual shape \texttt{nP} use
$m=(0,2,1,3)$, and for \texttt{nM} use $(0,1,2,3)$.
The residual representation becomes diagonal and every
flag becomes the identity. At a diagonal representation
every global coordinate of a tangent is a cocycle of the
corresponding root character (diagonal coordinates have
trivial character). Absent roots have zero cocycle.
Lower coordinates are locally coboundaries by (CP3).
In the rows of (CP2) satisfying either of the stated
support conditions there is at most one allowed upper
odd root among $A_0,B_0,H$ (the odd roots are all roots
except $\pm P$). Hence, after setting that single
coordinate to zero if necessary, all odd coordinates
are strict cocycles in the source's \texttt{Str} spaces.

In the remaining \texttt{Good} alternatives the residual
support is \texttt{vV} and its flag has a nonzero $(2,0)$
entry $b_v$ with support \texttt{qQ}. The degeneration
$m=(0,1,0,1)$ keeps just the flag
\[
                      1+b_v(E_{20}+E_{31})
\]
and makes the representation diagonal. The residual
local characters satisfy $\chi_0=\chi_2,\chi_1=\chi_3$:
in the original flag identity its $(2,0)$ entry gives
$b_v\chi_2=b_v\chi_0$ (row two has no off-diagonal
entry in support \texttt{vV}), and the paired product
gives the other equality. Write
$W=Z(s)+\operatorname{Ad}(\operatorname{diag}(\chi(s)))Y_v-Y_v$.
Direct multiplication of the limiting flag in (CP3),
using the symplectic transpose relation on $W$, gives
\[
 W_{10}=-b_v W_{12},\qquad
 W_{30}=-b_v^2W_{12},\qquad
 W_{21}=2b_vW_{01}+b_v^2W_{03}.                       \tag{CP4}
\]
On this decomposition group the three root characters
in the first two equations agree; those in the last
equation also agree. Modulo local coboundaries the
same equations consequently hold for the global
cocycle coordinates. Each coefficient displayed in
(CP4) is a unit, including $2b_v$. Thus if any of
the three negative coordinates is zero or strict, all
three are strict; if any two of the three positive
coordinates are zero or strict, the third is strict.

In this case $S$ must contain $-A_0$ or both $B_0,-H$.
The permitted rows of (CP2) are exactly
$1,3,5,6,7,9$. In rows $1,3,7$ set the odd coordinate
$-A_0$ to zero; in rows $5,6$ use $A_0$; in row $9$
use $H$. In every row after this one deletion the
negative triple $\{-A_0,B_0,-H\}$ is not fully present
and at most one of the positive triple
$\{A_0,-B_0,H\}$ is present. Formula (CP4) now gives
strictness of all odd coordinates. This argument
includes the coefficient prime three.

In either degeneration the one exceptional odd
coordinate has global cocycle dimension at most
$d+C$ by \texttt{Bounds.odd}. Its kernel maps
injectively to the ten independent symplectic
global matrix coordinates and the strictly lower
flag coordinates at all ordinary places. Each
global coordinate now belongs to a space of
dimension at most $C$, by \texttt{Bounds.str} for
odd roots and \texttt{Bounds.even} for the others.
The flag coordinates have dimension at most
$6r\leq6C$ (in fact four per symplectic flag).
The total tangent dimension is at most
$d+C+10C+6C\leq d+20C$. A minimal presentation
over the DVR then proves the required dimension
bound for the closed locus. No bound on the full
unrestricted adjoint tangent has been used.
\end{proof}

\begin{lemma}[Irreducible fibre bound]
\label{lem:rh-four-fibre}
Let $\mathfrak q$ be a prime of the completed chart $R$
at which the representation is absolutely irreducible.
Then
\[
 \dim R/\mathfrak q
 \leq\dim B/(\mathfrak q\cap B)+2+4r.                 \tag{CP5}
\]
\end{lemma}
\begin{proof}
Contract the prime to $\mathfrak p\subset A$. Put
$E=\operatorname{Frac}(B/(\mathfrak p\cap B))$ and
$E'=\operatorname{Frac}(A/\mathfrak p)$.
The marked diagonal eigenvalues $\lambda_i$ are
algebraic over $E$ and distinct. Adjoin them to
$E$ inside $E'$, obtaining $E_1$. Each coordinate
spectral projector is a polynomial in the marked
matrix with coefficients in $E_1$. Consequently
traces of words in these projectors and the
representation matrices are in $E_1$: on expanding
the projectors they become linear combinations of
traces of group elements, which belong to $B$.
In particular products of entries around every
directed cycle lie in $E_1$.

The directed graph of nonzero matrix entries
(row $i$, column $j$ denotes an edge $j\to i$)
is strongly connected, since otherwise a proper
nonempty terminal subset would span an invariant
coordinate subspace. Choose paths from vertex zero
to every vertex $i$, let $b_i$ be their nonzero
entry products and $b_0=1$. Choose also nonzero
return path products $c_i$. For every nonzero
entry $R(g)_{ij}$, the two cycle products
$b_j R(g)_{ij}c_i$ and $b_ic_i$ lie in $E_1$.
Hence the conjugation by
$\operatorname{diag}(b_0,b_1,b_2,b_3)$ makes every
representation entry lie in $E_1$.

The transformed alternating form has paired entries
proportional to $b_0b_3$ and $b_1b_2$. In the
transformed coordinates the equations for an
alternating form of the prescribed multiplier
have coefficients in $E_1$. Their solution space
over $E'$ is one-dimensional by absolute
irreducibility and Schur's lemma, using the
invertible transformed form. Scalar extension
preserves the dimension of this linear solution
space (one can take a finite spanning subset of
its defining equations). It follows that the
ratio $b_1b_2/b_3$ belongs to $E_1$.
All original representation entries therefore
belong to $E_1(b_1,b_2)$. Each lower symplectic
unipotent flag is determined polynomially by its
four coordinates $(1,0),(2,0),(3,0),(2,1)$:
its other two lower entries are
$-(1,0)$ and $(2,0)-(2,1)(1,0)$.
The original entry-generation property gives
$\operatorname{trdeg}_E E'\leq2+4r$.

Finally the local dimension inequality for a
finite type algebra over the complete Noetherian
local ring $B/(\mathfrak p\cap B)$ gives the same
upper bound for the dimension of
$(A/\mathfrak p)_{\ker x}$. Completion preserves
that local dimension, and $R/\mathfrak q$ is a
quotient of this completion. This proves (CP5);
it does not assert that a completed fraction field
has finite transcendence degree.
\end{proof}

\begin{theorem}[Component-support propagation and the residual chain]
\label{thm:rh-four-component-chain}
For the original fitted deep tower satisfying the stated
\texttt{Bounds} and degree threshold, the ideal $J$
satisfies the original local assertion
\[
 \texttt{SomeIn}(J,\ker x)\ \Longrightarrow\
 \texttt{AllIn}(J,\ker x)
\]
at every \texttt{IsRes}, \texttt{Good} residue pair.
In particular, starting at a seed, the same ideal
vanishes on every component through each endpoint
of the original \texttt{Chain}, and hence on
every lift of its terminal residue pair.
\end{theorem}
\begin{proof}
Let $N=2d+1-2r$. Grothendieck's connectedness theorem
for a complete regular local ring modulo $b$ equations,
applied to Lemma~\ref{lem:rh-four-chart-deficit}, makes
the graph of minimal primes of $R$ connected when
an edge joins $\mathfrak p,\mathfrak p'$ with
$\dim R/(\mathfrak p+\mathfrak p')\geq N+1$.
(The connectedness dimension is at least
$(a+1)-b-1$.) Suppose $\mathfrak p$ contains the
image of $J$. Choose a prime above their sum and
the coefficient uniformizer with quotient dimension
at least $N$, by the principal ideal theorem on
a maximal-dimensional component of the intersection.
All primes considered below contain $J$.

The threshold implies
\[
 d+20C+1<N,\qquad
 2d+1+2+4r<N+d_v.                                  \tag{CP6}
\]
The second inequality uses $d_v=d/r$ and $r\leq C$.
Thus the chosen prime cannot lie in any reducible
locus, by Lemma~\ref{lem:rh-four-red-dimension}.
Nor can it lie in any of the finitely many local
collision loci defined by equality of two labeled
characters on $T.U1(v)$. Indeed the earlier
weight-collision estimate bounds the dimension of
its coefficient quotient by $2d+1-d_v$, and
Lemma~\ref{lem:rh-four-fibre} adds at most $2+4r$.
This contradicts (CP6). This argument uses the
finite-weight property of the very same $J$.

One can now cut down to a one-dimensional prime
quotient avoiding all these bad loci. In the
complete local domain given by the chosen prime,
take a nonzero product of one witness outside
each bad ideal. In dimension greater than one,
choose a nonzero element of the maximal ideal
outside the finitely many height-one primes
containing that product. Every prime minimal
over the chosen principal ideal has height one
and avoids the product; the dimension drops by
one, since the local domain is complete and
catenary. Repeat until dimension one. The resulting
complete local domain is of characteristic $\ell$
and has finite residue field. Its normalization
is a finite extension, a complete DVR
$k'[[z]]$ for some finite field $k'$. This gives
an actual map $A\to k'[[z]]$ with the prescribed
residue on $B$, on which $J$ vanishes, the generic
representation is absolutely irreducible, and
each pair of ordinary characters is separated
on the required group. Theorem~\ref{thm:rh-four-p1}
therefore applies. Its all-primes-below-curve
conclusion, applied to the contraction of
$\mathfrak p'$, puts $J$ on that component as
well.

Iterating over the connected graph propagates
support to all chart minimal primes from any
one supported chart minimal prime. Flatness
$A\to R$ lifts each minimal prime of $A$ below
$\ker x$ to a minimal prime of $R$: local
completion is faithfully flat over $A_{\ker x}$,
and going-down gives the minimal-prime
assertion. Thus \texttt{SomeIn} implies
\texttt{AllIn} in the literal original space.

For a seed point, the previously proved
\texttt{LocZero} gives support on every minimal
prime of $A$ through its kernel, by contracting
to $B$. Each of these lies below the residue
kernel of its integral lift, so there is
\texttt{SomeIn} at the seed residue evaluation.
The residue of the seed is \texttt{Good} in the
original chain construction; hence it has
\texttt{AllIn}. An actual \texttt{Linked}
polynomial pair gives, by \texttt{Space.line},
a map $A\to k[t]$ inducing both endpoint
evaluations. Its prime kernel lies below
both endpoint kernels. A minimal prime
under this kernel is supported by
\texttt{AllIn} at the known endpoint. It
gives \texttt{SomeIn} at the other endpoint
and hence \texttt{AllIn} there whenever
that endpoint is \texttt{Good}. The
argument is symmetric in the endpoints.
Induction along the actual finite chain
establishes the assertion. For any map
from $A$ to a local domain inducing the
terminal residue evaluation, its prime
kernel contains a minimal prime through
that residue evaluation, so the image of
$J$ is zero.
\end{proof}

This is the component-support assertion of
\texttt{residual\_charpoly\_rank\_four}, with its original
arithmetic bounds and original degree threshold, together
with its use along the constructed residual links.
It proves support for local-domain points; no claim that
the same ideal annihilates nilpotent coefficient directions
is made here. The independent task of producing the
bounded-cocycle constants in the required arbitrarily deep
towers is retained. The strict local-window and full
arbitrary-quadratic descent clauses of \texttt{PMExAlt.aut}
also require a separate argument.
